\chapter{Suites arithmétiques, géométriques et sommes}
\section*{Exercice 5.1 -- Suite arithmétique}
On a $u_8-u_3=5r=15$, donc $r=3$ et $u_0=11-9=2$. Il y a $9$ termes de $u_0$ à $u_8$ ; la somme vaut $9(2+26)/2=126$.

\section*{Exercice 5.2 -- Suite géométrique}
La somme géométrique donne $5[1-(0{,}8)^{n+1}]/(1-0{,}8)=25[1-(0{,}8)^{n+1}]$. Puisque $(0{,}8)^{n+1}\to0$, la somme tend vers $25$. Chaque somme finie est strictement inférieure à $25$.

\section*{Exercice 5.3 -- Une raison négative}
La liste est $1,-2,4,-8,16$ ; la somme vaut $11$, aussi donnée par $(1-(-2)^5)/(1-(-2))=33/3$. Les différences changent de signe : la suite n'est ni croissante ni décroissante. Une raison inférieure à $1$ ne suffit pas à conclure à la décroissance.

\section*{Exercice 5.4 -- Somme décalée}
Les indices entiers de $4$ à $12$ inclus donnent $12-4+1=9$ termes. La suite est arithmétique, de premier terme $11$ et dernier terme $35$. La somme vaut $9(11+35)/2=207$.

\section*{Exercice 5.5 -- Télescopage}
La somme du membre de gauche de l'identité $(k+1)^3-k^3=3k^2+3k+1$ vaut $(n+1)^3-1$. Donc $3\sum k^2=(n+1)^3-1-3n(n+1)/2-n$. Développer et factoriser donne $n(n+1)(2n+1)/2$ ; diviser par $3$ termine.

\section*{Exercice 5.6 -- Somme des cubes}
On a $T_k^2-T_{k-1}^2=(T_k-T_{k-1})(T_k+T_{k-1})=k\cdot k^2=k^3$. La somme télescope donc en $T_n^2-T_0^2=T_n^2$. C'est un exemple de choix d'une bonne différence pour calculer une somme.
